% critique.tex — f169r, vv. 551–592 : commentaire critique
% Articles textuels tirés du commentaire de l'auteur (latex/Aristophane-Thesmo Harmattan 2023 vv 1-641.tex),
% texte et notes conservés ; répartition détaillée dans repartition.md.
\cc{556–7}{ὁρᾷς ὡς τλεγγίδας λαβοῦσαι / ἔπειτα σιφωνίζομεν τὸν σῖτον.}{Voilà un trait sur les mœurs féminines qui n'a pas été compris des différents exégètes ; c'est donc à moi que revient la tâche – ô combien délicate ! – d'expliquer ce qui répugne aux oreilles. Muses, venez m'aider... mais vous êtes pucelles ; Muses, ne bougez donc. Ah ! que l'on se sent seul ; je ne peux même pas lancer en avant l'honorable professeur Jeffrey Henderson, le plus fin limier des obscénités d'Aristophane et dont le groin fureteur parvient même à les débusquer là où elles ne sont pas (celle-ci, du moins, lui aura échappé.) Mais la recherche est un plaisir solitaire. Attelons-nous y donc ! Pris littéralement, le texte a de quoi surprendre : « tu vois comme nous saisissons les strigiles, et qu'ensuite nous siphonnons le pain. (sic) » Si vous êtes pris de l'envie de corriger, avec Küster, τὸν σῖτον en τὸν οἶνον (l'on siphonnera plutôt du vin que du pain), renseignez-vous sur ce qu'est un strigile (une sorte de spatule ou plaque rectangulaire tout en longueur, légèrement recourbée pour servir de racloir lors de la toilette) ; j'offre l'édition d'Aristophane de Küster à qui parviendra à se saouler en buvant avec ce genre de coupe. (Il faudra d'ailleurs un jour que l'on m'explique l'intérêt d'introduire des émendations si c'est pour laisser un texte tout aussi obscur – et retouché en prime.) Qu'à cela ne tienne, ce n'est pas ce genre de peccadille qui doit arrêter Austin \& Olson : Aristote a dit un jour (\italique{Topiques}, 145a24) qu'on pouvait se servir d'un strigile pour puiser de l'eau (hé ! ce sont des choses qui arrivent lorsqu'on est au bain ; A.\&O. auraient pu citer la phrase suivante qui leur était adressée~: Ἀλλ᾽ ὅμως εἴ τις ὁρίσαιτο τὴν στλεγγίδα ὄργανον πρὸς τὸ ἀρύειν, ἡμάρτηκεν, « Mais cependant, celui qui définirait le strigile comme instrument pour puiser se trompe »). Non, l'hypothèse d'ivrognesses s'enivrant à grandes rasades de strigiles est à écarter, sauf à m'expliquer : 1° comment boire au strigile \ \ \ 2° pourquoi aller chercher dans la salle de bain un instrument qui n'est pas censé en bouger \ \ \ 3° et quand bien même, qu'y aurait-il de déshonorant dans cette pratique ? \ \ \ Il faut chercher ailleurs, car les tentations des bains ne sont pas d'ordre biberonnesque. Aristophane vitupère ailleurs\nc{C'est un point de la querelle du Juste et de l'Injuste des \italique{Nuées} (v. 1044-1046) :
\citComm{Ἄδικος Λόγος — σε θερμῷ φησι λοῦσθαι πρῶτον οὐκ ἐάσειν.\\
καίτοι τίνα γνώμην ἔχων ψέγεις τὰ θερμὰ λουτρά;\\
Δίκαιος Λόγος — ὁτιὴ κάκιστόν ἐστι καὶ δειλὸν ποιεῖ τὸν ἄνδρα.\\
\francais{Raisonnement Injuste — il ne te permettra pas, dit-il, de te baigner dans l'eau chaude.\\
(Au Juste :) Mais d'après quel principe improuves-tu les bains chauds? \\
Raisonnement Juste — Parce qu'ils sont très pernicieux et qu'ils rendent l'homme lâche.}}Platon, d'accord sur ce point avec Aristophane comme sur beaucoup d'autres, nomme les bains chauds γεροντικὰ λουτρά, « des bains de vieillards » (\italique{Lois}, VI, 761c). Érasme en a fait un adage (2061): ἁπαλοὶ θερμολουσίαις, \italique{amollis par les bains chauds}.} contre cette propension des Athéniens décadents à se délicater : les bains chauds sont un exemple parmi d'autres de ces mauvaises habitudes qui amollissent les cœurs et durcissent les vices. Pour comprendre exactement où il veut en venir, c'est Euripide lui-même qui nous donnera la clef de la salle de bain : dans son drame satyrique du \italique{Cyclope}, le héros aux-mille-tours-lubriques, se désole sur « son cher siphon » (v. 339, τὸν σίφωνα τὸν φίλον) qui depuis trop longtemps vit dans la solitude – solitude vis-à-vis du beau sexe ; le verbe χηρεύω servant aussi à désigner l'état de veuvage. Il parle évidemment de son πέος~; une fois attesté ce sens pour σίφων, celui du verbe σιφωνίζω s'éclaire d'une lumière nouvelle (le verbe est d'autant plus amusant qu'il joue sur le rapprochement, cher à Aristophane, entre βινεῖν et πίνειν).\\
La présence étonnante de τὸν σῖτον s'explique alors aisément : la leçon originale était τὸν κύσθον $=$ \latin{pudenda muliebria}, très proche phonétiquement de τὸν σῖτον. 
Iotacisme aidant, il s'agit d'une simple \italique{métathèse} du `s' (avec chute du `k').
Il est remarquable de noter que la même métathèse se produit au vers 1114 : le rustre archer Scythe, dans l'optique de prouver que Mnésiloque n'est pas une femme, conseille d'aller y voir de plus près :\citComm{σκέψαι †τὸ σκύτο†.}Scaliger \italique{à la grâce du mètre} -- et surtout à la grâce du sens -- avait rétabli :\citComm{σκέψαι τὸ κύστο ($=$ τὸν κύσθον dans le parler scythe).}Il est à croire qu'un scribe des \italique{Thesmophories} confondait avec un autre le mot grec pour désigner le « plus bel apanage » du corps féminin, qui tentait tendrement tant Brassens \& Villon : il s'y sera trompé par deux fois dans notre pièce (\latin{bis repetita...}).
J'arrête mes explications ici, de peur qu'un Calliclès de passage ne m'assène :\citComm{Καλλίκλης-Φιλόλογος --- οὐκ αἰσχύνῃ εἰς τοιαῦτα ἄγων, ὦ Πέπωκες, τοὺς λόγους;\nc{Platon, \italique{Gorgias} 494e.}}Pour la suite (la coupe d'Onésimos, et l'explication d'un ami qui garde τὸν σῖτον), voir le commentaire de la pièce.}
